\section{Conclusiones}
El proposito de esta práctica fue evaluar los resultados obtenidos por un Agente Spice construido con los conocimientos vistos en clase. Adicionalmente el agente AI añadió conocimientos y scripts necesarios para el desarrollo y análisis de esta práctica. Los resultados se compararon con los análisis realizados usando HSPICE, dichos resultados se encuentran en las secciones 4, 5 y 6. 

Como se pudo notar por la evaluación puntual de algunas frecuencias mostradas en la comparación, los valores coinciden con un margen de error muy pequeño en el orden de $x10^{-7}$ esto  nos quiere decir que la similitud entre los valores esperados por parte de la simulación de HSPICE y los valores obtenidos por el agente son muy similares.
Adicionalmente se puede ver visualmente que las gráficas coinciden en comportamiento, siendo muy similares entre si.
Algunas carácteristicas adicionales que proporciona este agente es la obtención del grafo del circuito y matríz de Tableua que representa matemáticamente el circuito analizado.

Por lo que se puede concluir que si se logró realizar un analizador de circuitos tipo SPICE usando un Agente AI en python. Se obtuvieron resulltados muy similares a los de la simulación de HSPICE. Algunas deficiencias que tiene este Agente SPICE sería para simular circuitos más complejos que incluyan transistores y que requieran modelados más complejos.
